\include{settings}

% 填写个人信息
% 学院 / 部门
\newcommand{\school}{人工智能学域 | AI Thrust} 

% 联系方式（调整样式：去掉白色文字，适配个人信息板块）
\newcommand{\contact}{
    \scriptsize
    % 邮箱 
    \faEnvelope \quad \href{mailto:youremail@hust.edu.cn}{xxxxxxxxxx}
    \hspace{2em}
    % 手机号
    \faPhone \quad  111-1111-1111
    \hspace{2em}
    % 微信/Github
    \faGithub \quad \href{https://github.com/yourname}{https:xxxxxx}
}

%%%%%%%%%%%%%%%%%%%%
% 简历正文
%%%%%%%%%%%%%%%%%%%%
\begin{document}

    \begin{tikzpicture}[remember picture, overlay]
        % 1. 画背景条
        \fill[HKUST_GZ] (current page.north west) rectangle ([yshift=-1.6cm]current page.north east);
        
        % 2. 放置左侧校徽文字 Logo 
        % 如果没有白字Logo，可以将下面这行注释掉，直接用文字写校名
        \node[anchor=west] at ([xshift=1cm, yshift=-0.8cm]current page.north west){
            \includegraphics[height=1.0cm]{images/UST-GZ_white_tc_L_PNG.png} 
        };
        % 3. 放置右侧学院名称
        \node[anchor=east] at ([xshift=-1cm, yshift=-0.8cm]current page.north east){
            \textcolor{white}{\textbf{\school}}
        };
    \end{tikzpicture}
    \vspace{-2em} % 调整正文与页眉的距离


    % ============================================================
    % 背景水印 
    % ============================================================
    \begin{tikzpicture}[remember picture, overlay]
        \node[
            opacity=0.1,           % 透明度：建议 0.1 到 0.15 之间 (太深会干扰文字)
            blend mode=multiply,    % 混合模式：让图片的白色背景变透明 (类似PS的正片叠底)
            inner sep=0pt           % 去除多余边距
        ] at(current page.center){
            \includegraphics[width=0.8\paperwidth, keepaspectratio]{images/in_UST-GZ_4C_en-C_PNG.png}
        };
    \end{tikzpicture}

    % ============================================================
    % 个人信息板块（新增联系方式）
    % ============================================================
    \begin{figure}[h]
        % 左半边：文字信息
        \begin{minipage}{0.82\textwidth}
            \section{\makebox[\widthof{\faAddressCard}][c]{\color{HKUST_GZ}{\faAddressCard}}\quad 个人信息}
            \begin{tabularx}{\linewidth}{p{\widthof{出生日期:}}Xp{\widthof{政治面貌:}}X}
                姓\ \ \ \ \ \ \ \ 名: & XXXX & 
                性\ \ \ \ \ \ \ \ 别: & XX  \\
                年\ \ \ \ \ \ \ \ 龄: & XX & 
                政治面貌: & XXXX \\
                联系方式: & \multicolumn{3}{l}{\contact} \\
            \end{tabularx}
        \end{minipage}
        \hspace{1em}
        % 右半边：照片
        \begin{minipage}{0.14\textwidth}
            \setlength{\fboxsep}{0pt}
            \fcolorbox{HKUST_GZ}{white}{\includegraphics[width=\linewidth]{images/avatar.png}}
        \end{minipage}
    \end{figure}
    \vspace{-1.5em}

    % ============================================================
    % 教育背景
    % ============================================================
    \section{\makebox[\widthof{\faGraduationCap}][c]{\color{HKUST_GZ}{\faGraduationCap}}\quad 教育背景}

    % 示例内容
    {\large \textbf{XXXXXX}}，本科 \hfill {xxxx，xxxx} \\
    {{XXXXX}}，专业：XXXX \hfill {XXXX-XXXX} \\

    \vspace{0.5em}
    {\large \textbf{香港科技大学(广州)}}，硕士 \hfill {广东，广州} \\
    {{信息枢纽}}，专业：人工智能 \hfill {2025年9月-至今} \\
    \textbf{研究方向}：XXXXXXXXXXX。

    % ============================================================
    % 科研成果
    % ============================================================
    \section{\makebox[\widthof{\faBook}][c]{\color{HKUST_GZ}{\faBook}}\quad 科研成果}
     \vspace{0.5em}
    \textbf{XXXXXXXXXX} \\
    \textbf{第一作者} \hfill 
    发表于 \textbf{XXXXXX 2026} (CCF-A类会议)

    
    % ============================================================
    % 项目与实习
    % ============================================================
    \section{\makebox[\widthof{\faLaptopCode}][c]{\color{HKUST_GZ}{\faLaptopCode}}\quad 项目与实习经历}
    
    \vspace{0.5em}
    {\large{\textbf{xxxxx}}}，大模型算法实习生 \hfill {2025年4月 - 2025年9月}\\
    \textbf{主要职责}：负责多模态大模型的微调与评测。\\
    \textbf{工作产出}：xxxxx。

% \newpage
%换页则复制下面的代码
    % \begin{tikzpicture}[remember picture, overlay]
    %     \node[
    %         opacity=0.1,           % 透明度：建议 0.1 到 0.15 之间 (太深会干扰文字)
    %         blend mode=multiply,    % 混合模式：让图片的白色背景变透明 (类似PS的正片叠底)
    %         inner sep=0pt           % 去除多余边距
    %     ] at(current page.center){
    %         \includegraphics[width=0.8\paperwidth, keepaspectratio]{images/in_UST-GZ_4C_en-C_PNG.png}
    %     };
    % \end{tikzpicture}


    % ============================================================
    % 竞赛
    % ============================================================
        \section{\makebox[\widthof{\faWrench}][c]{\color{HKUST_GZ}{\faWrench}}\quad 竞赛}
    \vspace{0.5em}
    {\large{\textbf{xxxxx}}}，xxxxx \hfill {xxxxx}\\
    \textbf{主要职责}：xxxxx\\
    \textbf{工作产出}：xxxxx
   

    
    % ============================================================
    % 奖项
    % ============================================================
        \section{\makebox[\widthof{    \faTrophy
}][c]{\color{HKUST_GZ}{\faTrophy}}\quad 奖项}
    \vspace{0.5em}
    {\large{\textbf{xxxxx}}} \hfill {xxxxx}\\

\end{document}